\scanpage{118}

Cette condition n'a pas l'air tautologique !

D'autre part, la condition de connexité pour $A_\lambda(s)$ peut s'examiner à l'aide de la bijection canonique (36) induite par $\lambda$

\[
A_0(s)\xrightarrow{\ \sim\ } A_\lambda(s)
\]

on note que $\{\lambda t,\lambda t'\}\in A_\lambda$ \unsure{sss|ssi} $\{t,\lambda t'\}\in A_0$ -- -- pourquoi cette relation \inserted{symétrique} entre deux éléments $t,t'$ de $A_0(s)$ serait elle ``connexe'' ?

% slanted
\margin{on voit bien, sur l'exemple de la double pyramide (sur un $n$-polygone avec $n$ pair $\geqslant6$, $n\not\equiv0\ (4)$) que cette condition n'est \uline{pas} tautologique !}

En fait, on va \unsure{apaiser|applanir} ces perplexités par un axiome qui, contrairement à l'axiome ``par acquit de \unsure{conscient|conscience}'' $T_0\,4$, a l'air des plus substantiels

\[
\begin{aligned}
(T_0\,5)\qquad &\forall\lambda\in Z',\quad\text{la structure de graphe}\\
&\Gamma_\lambda=(S,A_\lambda)\quad(\text{cf }35)\quad\text{est \uline{isomorphe}}\\
&\text{à }\Gamma_0=(S,A_0),
\end{aligned}\tag*{$(41)$}
\]

en d'autres termes, on demande qu'il existe une permutation $g=g_\lambda\in\unsure{\mathfrak S_S|G_S}$ telle que l'on ait

\[
\forall s,t\in S,\ \text{on a}\quad \{s,t\}\in A_0\Longleftrightarrow
 \underbrace{\{gs,\lambda^{-1}gt\}\in A_0}_{\substack{\text{ce qui exprime}\\\{gs,gt\}\in A_\lambda}}
\tag*{$(42)$}
\]
% archive sheet number: 117
